
\documentclass[11pt]{article}

% Page and typography
\usepackage[a4paper,margin=25mm,headheight=15pt]{geometry}
\usepackage[T1]{fontenc}
\usepackage[utf8]{inputenc}
\usepackage{lmodern}
\usepackage{microtype}
\usepackage{setspace}

% Mathematics, tables, and figures
\usepackage{amsmath,amssymb}
\usepackage{graphicx}
\usepackage{booktabs}
\usepackage{tabularx}
\usepackage{array}
\usepackage{caption}
\usepackage{float}
\usepackage[section]{placeins}
\usepackage{needspace}

% Document styling
\usepackage[dvipsnames]{xcolor}
\usepackage{fancyhdr}
\usepackage{titlesec}
\usepackage{enumitem}
\usepackage{hyperref}

\definecolor{reportblue}{HTML}{17365D}
\definecolor{reportgray}{HTML}{5B6573}
\definecolor{rulegray}{HTML}{C9CED6}

\hypersetup{
  colorlinks=true,
  linkcolor=reportblue,
  citecolor=reportblue,
  urlcolor=reportblue,
  pdftitle={Spatial Patterns in Indian District Night-Light Growth},
  pdfauthor={Nakshatra Ghosh},
  pdfsubject={Descriptive statistics, spatial econometrics, and cluster analysis}
}

\setlength{\parindent}{0pt}
\setlength{\parskip}{0.55\baselineskip}
\setstretch{1.04}
\setcounter{tocdepth}{2}
\setlist[itemize]{leftmargin=1.5em,itemsep=0.2em,topsep=0.25em}

\titleformat{\section}
  {\Large\bfseries\color{reportblue}}
  {\thesection}{0.65em}{}
  [\vspace{-0.35em}\color{rulegray}\titlerule]
\titleformat{\subsection}
  {\large\bfseries\color{reportblue}}
  {\thesubsection}{0.6em}{}

\captionsetup{
  font=small,
  labelfont={bf,color=reportblue},
  justification=raggedright,
  singlelinecheck=false,
  skip=6pt
}

\pagestyle{fancy}
\fancyhf{}
\fancyhead[L]{\small\color{reportgray}Indian district night-light growth}
\fancyhead[R]{\small\color{reportgray}SHRUG v2.2}
\fancyfoot[C]{\small\thepage}
\renewcommand{\headrulewidth}{0.4pt}
\renewcommand{\footrulewidth}{0pt}

\newcommand{\V}[1]{\mathbf{#1}}
\newcommand{\W}{\mathbf{W}}
\newcommand{\Rtwo}{\ensuremath{R^{2}}}
\newcommand{\NA}{\textemdash}
\newcommand{\reportsection}[1]{\Needspace{8\baselineskip}\section{#1}}
\newcolumntype{Y}{>{\raggedright\arraybackslash}X}
\newcolumntype{C}{>{\centering\arraybackslash}X}

\begin{document}

\hypersetup{pageanchor=false}
\begin{titlepage}
  \thispagestyle{empty}
  \vspace*{18mm}

  {\color{reportblue}\rule{\textwidth}{2.2pt}}\\[12mm]
  {\fontsize{27}{32}\selectfont\bfseries\color{reportblue}
    Spatial Patterns in Indian District\\[2mm]
    Night-Light Growth\par}
  \vspace{7mm}
  {\Large\color{reportgray}
    Descriptive statistics, spatial econometrics, and cluster analysis\par}
  \vspace{11mm}
  {\color{rulegray}\rule{\textwidth}{0.8pt}}

  \vfill

  \begin{tabular}{@{}ll@{}}
    \textbf{Author} & Nakshatra Ghosh \\
    \textbf{Document} & Research report \\
    \textbf{Data source} & SHRUG v2.2 \\
    \textbf{Date} & 30 September 2026
  \end{tabular}

  \vfill

  \begin{minipage}{0.92\textwidth}
    \small\color{reportgray}
    This report studies growth in satellite-observed nighttime light across
    640 Indian Census 2011 districts. Results describe associations in
    remotely sensed light; they are not estimates of district GDP or causal
    spatial spillovers.
  \end{minipage}
  \vspace{10mm}
\end{titlepage}

\hypersetup{pageanchor=true}
\pagenumbering{roman}

\section*{Abstract}
\addcontentsline{toc}{section}{Abstract}

This study analyzes 640 Indian Census 2011 districts using SHRUG v2.2
nighttime lights. The main comparison uses VIIRS observations from 2013 to
2021 and fixed 2011 population to normalize district size. Lower initially lit
districts show faster subsequent light growth in a conditional model, while
growth is strongly spatially clustered. Spatial lag, spatial error, and
spatial Durbin models separate forms of dependence. Local Moran analysis and
a graph-constrained Ward typology describe regional clusters. These are
associations in satellite light, not district GDP or causal spillovers.

\subsection*{Key results}

\begin{table}[H]
  \centering
  \caption*{Headline estimates}
  \begin{tabularx}{\textwidth}{@{}Y>{\raggedleft\arraybackslash}p{0.31\textwidth}@{}}
    \toprule
    \textbf{Measure} & \textbf{Estimate} \\
    \midrule
    Mean annual 2013--2021 log-light growth & 0.0682 \\
    Conditional OLS initial-light coefficient & $-0.0457$ (HC3 SE 0.00344) \\
    Global Moran's $I$ for growth & 0.654 (permutation $p=0.001$) \\
    Spatial Durbin $\rho$ & 0.3253 \\
    Spatial Durbin residual Moran's $I$ & $-0.0037$ ($p=0.935$) \\
    FDR-significant high--high / low--low districts & 37 / 37 \\
    \bottomrule
  \end{tabularx}
\end{table}

\textbf{Interpretation.} Light is normalized by each district's 2011
population, held fixed in every year. It is therefore
baseline-population-normalized light, not contemporaneous per-capita income.

\clearpage
\tableofcontents
\clearpage

\pagenumbering{arabic}

\reportsection{Research question and data}

The first project in the supplied brief asks whether less initially lit Indian
districts grew faster, whether growth forms spatial clusters, and whether
nearby outcomes or characteristics are associated with district growth. The
unit is the Census 2011 district, held fixed over time. Night lights are an
imperfect economic-activity proxy: electrification, lighting technology, and
settlement patterns also matter \cite{henderson2012}.

The source tables are SHRUG v2.2 district VIIRS annual sum, calibrated DMSP
district light, the 2011 Population Census Abstract, a population/land-area
key, and district polygons \cite{shrug,asher2021}. The raw VIIRS table covers
2012--2023, but published metadata describes 2012--2021; the documented range
is used here. VIIRS and DMSP are never spliced.

\subsection{Data audit}

\begin{table}[H]
  \centering
  \caption{Data-integrity checks.}
  \begin{tabularx}{\textwidth}{@{}p{0.36\textwidth}Y@{}}
    \toprule
    \textbf{Check} & \textbf{Result} \\
    \midrule
    District census and area rows & 640 each; population values agree \\
    VIIRS derived panel & 12,800 rows $=640\times10$ years $\times2$ masks \\
    DMSP derived panel & 12,800 rows $=640\times20$ years \\
    Polygon match & 640 exact keys; one unlabeled district-000 placeholder removed \\
    Main endpoint sample & 640 districts, no missing covariates, positive endpoint VIIRS light \\
    \bottomrule
  \end{tabularx}
\end{table}

\subsection{Variable construction}

For district $i$ and year $t$, normalized light, log light, and annual growth
are defined as
\begin{align}
  L_{it} &= \frac{1{,}000\times \mathrm{VIIRS\_sum}_{it}}
                  {\mathrm{population}_{i,2011}}, \\
  z_{it} &= \log(L_{it}), \\
  g_i &= \frac{z_{i,2021}-z_{i,2013}}{8}.
\end{align}
The fixed denominator means that the growth rate equals growth of total
district light. Baseline controls are log 2011 density, adult literacy,
agricultural main-worker share, Scheduled Caste/Tribe share, and state
indicators.

The principal log comparison starts in 2013 because the 2012 VIIRS
observations include zeros. For DMSP, concurrent calibrated satellite
measurements are averaged within district-year. The historical log check
begins in 2000, the first year positive for every district.

\reportsection{Numerical and visual description}

\begin{table}[H]
  \centering
  \caption{District-level descriptive statistics.}
  \small
  \begin{tabularx}{\textwidth}{@{}Yrrrr@{}}
    \toprule
    \textbf{District variable} & \textbf{Mean} & \textbf{Median} & \textbf{P10} & \textbf{P90} \\
    \midrule
    VIIRS per 1,000 baseline residents, 2013 & 7.689 & 6.068 & 1.641 & 15.259 \\
    VIIRS per 1,000 baseline residents, 2021 & 11.209 & 9.829 & 4.653 & 18.272 \\
    Annual log-light growth & 0.0682 & 0.0588 & $-0.0004$ & 0.1478 \\
    2011 density & 994.0 & 409.8 & 137.7 & 1,251.8 \\
    2011 literacy share & 0.723 & 0.722 & 0.590 & 0.854 \\
    2011 agricultural worker share & 0.533 & 0.591 & 0.197 & 0.773 \\
    \bottomrule
  \end{tabularx}
\end{table}

Mean normalized light rises 45.8\% from 2013 to 2021; the median rises
62.0\%. Mean annual log growth is 0.0682, and 10.5\% of districts have
negative endpoint growth. The distribution is right-skewed. The annual series
has a noticeable 2016--2017 step, which could reflect illumination or
processing as well as activity.

\begin{figure}[H]
  \centering
  \includegraphics[width=0.94\textwidth]{results/figures/viirs_trend.png}
  \caption{VIIRS district mean and median light per 1,000 fixed 2011 residents.}
  \label{fig:viirs-trend}
\end{figure}

\clearpage
\reportsection{Distribution and spatial clustering}

\begin{figure}[H]
  \centering
  \includegraphics[width=\textwidth]{results/figures/distribution_convergence.png}
  \caption{Growth distribution and unconditional OLS line with a 95\% confidence band for the fitted mean.}
  \label{fig:distribution-convergence}
\end{figure}

The fitted unconditional line is
\[
  \widehat{g}=0.16662-0.05737\log(\text{initial normalized VIIRS light}),
  \qquad \Rtwo=0.586.
\]
The negative slope is an association in annual log-light growth, not a line
for GDP growth. The conditional fitted relation instead has an initial-light
slope of $-0.04569$ after controls and state effects (Section~\ref{sec:theory}).

Global Moran's $I$ compares district deviations with the weighted deviations
of Queen-contiguous neighbors. The Queen graph has 4.14 neighbors per district
on average, 34 disconnected components, and 12 districts with no contiguous
neighbor. Those 12 retain zero rows in $\W$; no distant boundary is invented.
An eight-nearest-neighbor graph is a sensitivity test.

\begin{table}[H]
  \centering
  \caption{Global spatial-autocorrelation tests.}
  \begin{tabularx}{0.86\textwidth}{@{}Yrr@{}}
    \toprule
    \textbf{Moran test} & \textbf{$I$} & \textbf{Two-sided permutation $p$} \\
    \midrule
    Initial normalized light & 0.6124 & 0.001 \\
    Annual growth & 0.6544 & 0.001 \\
    Conditional OLS residual & 0.1693 & 0.001 \\
    Spatial Durbin residual & $-0.0037$ & 0.935 \\
    \bottomrule
  \end{tabularx}
\end{table}

Local Moran statistics use 999 permutations and Benjamini--Hochberg false
discovery rate control at 5\% over the 628 districts with a Queen neighbor.
Results are 37 high--high, 37 low--low, two low--high, 552 not significant,
and 12 with no contiguous neighbor. High--high examples include districts in
Bihar and Manipur. Spatial clustering does not by itself establish diffusion.

\reportsection{Where growth clusters occur}

\begin{figure}[H]
  \centering
  \includegraphics[width=\textwidth]{results/figures/growth_maps.png}
  \caption{Annual 2013--2021 growth and Local Moran classes after FDR correction. Dark gray marks no-neighbor districts.}
  \label{fig:growth-maps}
\end{figure}

The map shows broad geographic heterogeneity, including fast-growing clusters
in parts of Bihar and the northeast. Neighbor definitions matter where the
source polygons have gaps or inaccuracies; the supplied SHRUG map
documentation cautions about boundary precision.

\subsection{Graph-constrained Ward grouping}

We standardize five district features: initial log light, annual growth, log
density, literacy, and agricultural-employment share. Ward agglomerative
clustering may merge districts only along a symmetric eight-nearest-neighbor
graph. We compare two through six groups with the silhouette score. Three
groups score highest (0.158), and each is connected in the graph. The low
score signals modest separation, so these are a descriptive typology, not
sharply defined economic regions. Growth is an input to clustering and cannot
be used as an independent validation.

\clearpage
\reportsection{Spatially constrained district typology}

\begin{figure}[H]
  \centering
  \includegraphics[width=0.92\textwidth,height=0.64\textheight,keepaspectratio]{results/figures/spatial_regions.png}
  \caption{Three Ward groups constrained by a symmetric eight-nearest-neighbor graph.}
  \label{fig:spatial-regions}
\end{figure}

\begin{table}[H]
  \centering
  \caption{Profiles of the three spatially constrained groups.}
  \begin{tabular}{@{}rrrrr@{}}
    \toprule
    \textbf{Group} & \textbf{Districts} & \textbf{Mean growth} &
    \textbf{Mean initial log light} & \textbf{Median density} \\
    \midrule
    1 & 252 & 0.1023 & 1.443 & 351.6 \\
    2 & 347 & 0.0406 & 1.999 & 495.5 \\
    3 & 41  & 0.0924 & 0.997 & 56.6 \\
    \bottomrule
  \end{tabular}
\end{table}

Group 1 combines relatively low initial light and fast growth; group 2 has
higher initial light and slower growth; group 3 has very low density and fast
growth. These profiles follow in part from the features used to form the
groups.

\reportsection{Econometric theory and implementation}
\label{sec:theory}

\subsection{OLS benchmark and convergence}

The absolute model regresses annual growth $g_i$ on initial log light
$z_{i,2013}$. The conditional model adds 2011 district characteristics and
state indicators. HC3 robust standard errors allow heteroskedasticity. A
negative coefficient is consistent with initially darker districts catching
up in light, conditional on measured covariates. Under a restrictive
transitional growth law,
\[
  \beta=\frac{\exp(-8\lambda)-1}{8},
  \qquad
  \lambda=-\frac{\log(1+8\beta)}{8}
\]
is an implied speed.

With state 01 as the reference, the estimated conditional equation is
\begin{align*}
  \widehat{g}_i={}&0.22351
  -0.04569\,\text{initial log light}_i \\
  &-0.01563\,\text{log density}_i
  +0.01361\,\text{literacy}_i \\
  &-0.01720\,\text{agricultural worker share}_i
  -0.01023\,\text{SC/ST share}_i
  +\widehat{\text{state effect}}_i.
\end{align*}
The intercept is a model reference term; the plotted line in
Figure~\ref{fig:distribution-convergence} is the simpler unconditional fit.

\subsection{Three spatial alternatives}

Let $\V{Z}$ denote the design matrix and $\W$ the row-standardized spatial
weights matrix. The three alternatives are
\begin{align}
  \text{SAR:}\quad & \V{g}=\rho\W\V{g}+\V{Z}\V{b}+\V{\varepsilon}, \\
  \text{SEM:}\quad & \V{g}=\V{Z}\V{b}+\V{u},
    & \V{u}=\lambda\W\V{u}+\V{\varepsilon}, \\
  \text{SDM:}\quad & \V{g}=\rho\W\V{g}+\V{Z}\V{b}
    +\W\V{Z}\V{\theta}+\V{\varepsilon}.
\end{align}
The simultaneous spatial-outcome term in SAR requires maximum likelihood;
$\rho$ captures conditional spatial association. SEM attributes dependence to
unobserved spatial shocks. SDM additionally includes neighbors' initial light
and other continuous characteristics. State indicators are never spatially
lagged. All spatial models use Gaussian maximum likelihood with the Ord
eigenvalue method.

In the SDM, impacts require
\[
  \V{S}_k=(\V{I}-\rho\W)^{-1}(\beta_k\V{I}+\theta_k\W).
\]
The average direct impact is $\operatorname{tr}(\V{S}_k)/n$; the average total
impact is $\V{1}^{\mathsf T}\V{S}_k\V{1}/n$; and the indirect impact is total
minus direct. Exact matrix calculations retain the zero-neighbor rows.
Uncertainty intervals draw 5,000 times from the full model covariance matrix.
They do not cover uncertainty about $\W$ or the source measurement.

\subsection{Diagnostics}

We compare AIC, a nested likelihood-ratio test for SAR versus SDM, and Moran
tests on model residuals. OLS standard errors are HC3; spatial-model standard
errors use the Gaussian likelihood. Neither specification delivers causal
identification.

\reportsection{Model estimates and interpretation}

\begin{table}[H]
  \centering
  \caption{OLS and spatial-model estimates.}
  \small
  \begin{tabularx}{\textwidth}{@{}Yrrrr@{}}
    \toprule
    \textbf{Model} & \textbf{Initial coefficient} & \textbf{Spatial parameter} &
    \textbf{AIC} & \textbf{Residual $I$} \\
    \midrule
    OLS: initial level only & $-0.0574$ & \NA & \NA & \NA \\
    OLS: controls + state & $-0.0457$ & \NA & $-2637.6$ & 0.1693 \\
    SAR: controls + state & $-0.0433$ & $\rho=0.1966$ & $-2658.3$ & 0.0690 \\
    SEM: controls + state & $-0.0459$ & $\lambda=0.3169$ & $-2672.9$ & 0.0114\textsuperscript{*} \\
    SDM: controls + state & $-0.0468$ & $\rho=0.3253$ & $-2673.7$ & $-0.0037$ \\
    \bottomrule
  \end{tabularx}
  \vspace{0.35em}

  \parbox{\textwidth}{\footnotesize
  \textsuperscript{*}SEM reports Moran's $I$ for its filtered innovations.
  All models use 640 districts.}
\end{table}

The OLS conditional initial-light coefficient has HC3 SE 0.00344 and
$\Rtwo=0.800$. Its implied speed is 5.69\% per year (approximately
4.67\%--6.80\% from its coefficient interval), conditional on the
transitional-growth interpretation.

SAR leaves modest residual spatial dependence ($I=0.0690$, $p=0.024$). SEM
filtered innovations ($I=0.0114$, $p=0.665$) and SDM residuals
($I=-0.0037$, $p=0.935$) do not show detectable residual clustering. SDM has
the lowest AIC, but its 0.84-point advantage over SEM is small. The nested
SAR-versus-SDM likelihood-ratio statistic is 25.40 on five extra spatially
lagged regressors ($p=0.000117$). The data support additional spatial
structure compared with SAR without decisively separating covariate
spillovers from spatially correlated omitted factors.

\begin{table}[H]
  \centering
  \caption{Spatial Durbin impacts of initial log light.}
  \begin{tabular}{@{}lrrr@{}}
    \toprule
    \textbf{Impact} & \textbf{Point estimate} & \textbf{95\% lower} & \textbf{95\% upper} \\
    \midrule
    Average direct & $-0.04652$ & $-0.05089$ & $-0.04221$ \\
    Average indirect & 0.00428 & $-0.00358$ & 0.01222 \\
    Average total & $-0.04225$ & $-0.05107$ & $-0.03359$ \\
    \bottomrule
  \end{tabular}
\end{table}

The neighboring initial-light coefficient in the SDM is positive, but the
average indirect-impact interval crosses zero. It is therefore not evidence
for a precisely signed average spillover. The negative direct and total
impacts align with convergence in this light proxy. The outcome-spatial
parameter is an association, not proof that one district caused another to
grow.

\reportsection{Robustness checks}

\begin{table}[H]
  \centering
  \caption{Alternative samples, periods, measures, and spatial weights.}
  \footnotesize
  \begin{tabularx}{\textwidth}{@{}Yrrr@{}}
    \toprule
    \textbf{Specification} & \textbf{Districts} &
    \textbf{Conditional OLS $\beta$ (HC3 SE)} & \textbf{SAR $\rho$} \\
    \midrule
    VIIRS median-masked & 640 & $-0.0458$ (0.0035) & 0.214 \\
    VIIRS 2014--2019 & 640 & $-0.0536$ (0.0044) & 0.252 \\
    VIIRS eight-neighbor graph & 640 & $-0.0457$ (0.0034) & 0.141 \\
    Two-year-smoothed endpoints & 640 & $-0.0445$ (0.0036) & 0.238 \\
    Queen non-islands & 628 & $-0.0460$ (0.0035) & 0.184 \\
    Separate DMSP 2000--2013 & 640 & $-0.0259$ (0.0026) & 0.450 \\
    Trim 1\% in both growth tails & 626 & $-0.0421$ (0.0032) & \NA \\
    \bottomrule
  \end{tabularx}
\end{table}

The negative initial-light association persists under alternative VIIRS
masking, a pre-2020 interval, endpoint smoothing, geographic definitions,
removal of Queen islands, and trimming. The DMSP result is a historical sign
check only; its coefficients must not be compared numerically with VIIRS due
to sensor and processing differences.

\begin{figure}[H]
  \centering
  \includegraphics[width=0.94\textwidth]{results/figures/ols_residuals.png}
  \caption{Conditional OLS residuals versus fitted annual log growth.}
  \label{fig:ols-residuals}
\end{figure}

\reportsection{Limits and conclusion}

First, this is an analysis of lights, not directly observed district GDP or
consumption. Electrification, saturation, blooming, masking, and local
economic composition can alter the lights--output relation. Second, fixed
2011 population does not measure migration or population growth by 2021.
Third, state controls cannot remove district-specific omitted variables,
common regional shocks, or measurement error. The growth formula reuses the
initial light observation, so regression to the mean may contribute to its
negative coefficient. Fourth, the spatial models rely on a chosen neighbor
graph and Gaussian likelihood; geometry gaps produce disconnected Queen
components.

The supported conclusions are narrower: normalized light rose in most
districts; initially darker districts tended to show faster subsequent light
growth; and that growth was strongly spatially patterned. SEM and SDM explain
more of the residual geographic pattern than a conventional regression. A
stronger claim about economic convergence or causal spatial spillovers would
require validated output or consumption measures, annual population, and a
research design with exogenous sources of variation.

\subsection{Reproducibility and data rights}

The accompanying repository contains harmonized panels, district polygons,
numerical tables, figures, the Python pipeline, source links, and LaTeX
source. SHRUG-derived data and maps retain the upstream noncommercial CC
BY-NC-SA 4.0 terms. Data lineage and exact variable definitions appear in
\texttt{DATA\_SOURCES.md} in the full analysis repository.

\clearpage
\begin{thebibliography}{9}

\bibitem{asher2021}
Asher, S., Lunt, T., Matsuura, R., and Novosad, P. (2021).
Development research at high geographic resolution: An analysis of night
lights, firms, and poverty in India using the SHRUG open data platform.
\emph{World Bank Economic Review}, 35(4), 845--871.
\href{https://doi.org/10.1093/wber/lhab003}{doi:10.1093/wber/lhab003}.

\bibitem{henderson2012}
Henderson, J. V., Storeygard, A., and Weil, D. N. (2012).
Measuring economic growth from outer space.
\emph{American Economic Review}, 102(2), 994--1028.
\href{https://doi.org/10.1257/aer.102.2.994}{doi:10.1257/aer.102.2.994}.

\bibitem{shrug}
Development Data Lab. \emph{SHRUG v2.2: Socioeconomic High-resolution Rural--Urban Geographic Platform for India}.
Project and data documentation; license terms.

\bibitem{pysal}
PySAL Developers. \emph{spreg: Spatial Regression Models}.
Official documentation for spatial likelihood and impact estimation.

\end{thebibliography}

\end{document}
